%%% Local Variables:
%%% TeX-command-extra-options: "-shell-escape"
\documentclass{uksrc-memo}
\usepackage{microtype}
\usepackage{amsmath}
\usepackage{listings}
\usepackage[hidelinks]{hyperref}

\SetMemoMetadata{
  title={Example UKSRC Memo},
  id={UKSRC-MEMO-0001},
  type={MEMO},
  revision={0.1},
  % Allowed values: Draft | In Review | Approved | Superseded
  status={Draft},
  tags={validation, EoR, \JiraTag{SE-397}},
  classification={Unrestricted},
  authors={First Author, Second Author, Third Author},
  footer-authors={First Author et al.},
  contact={First Author (email@example.com)},
  % Choose one:
  % issue-date={07 April 2026},
  issue-date={\today},
  summary={This memo summarises the purpose, scope, main conclusions, and any actions requested of the reader. Replace this paragraph with a concise executive summary suitable for circulation as a stand-alone overview.}
}
\renewcommand{\uksrcdisclaimertext}{%
The UKSRC memos are intended to support the rapid recording of technical work, investigations, design notes, and open questions. A memo may reflect the views of its authors at the time of writing and should not automatically be taken as formal UKSRC policy, project baseline, or final design guidance.
}

\ShowAuthorApprovals
\ResetAuthorApprovals
\AddAuthorApprovalSection{Authored by}
\AddAuthorApproval{First Author}{Lead Author}{UKSRC}{[Signature] [Date]}
\AddAuthorApproval{Second Author}{Contributor}{Organisation}{}
\AddAuthorApprovalSection{Approved and released by}
\AddAuthorApproval{Reviewer Name}{Reviewer}{Organisation}{[Signature] [Date]}

\ResetDocumentHistory
\AddDocumentRevision{0.1}{\IssueDate}{Initial template draft}

\ResetAbbreviations
\AddAbbreviation{Abbreviation}{Description}
\AddAbbreviation{SKA}{Square Kilometre Array}
\AddAbbreviation{UKSRC}{United Kingdom SKA Regional Centre}
\AddAbbreviation{UKSRC-ST}{UKSRC Science Tooling}

\newcommand{\unit}[1]{\mathrm{#1}}
\newcommand{\unitp}[2]{\ensuremath{\mathrm{#1}^{#2}}}
\newcommand*{\fullref}[1]{\hyperref[{#1}]{\autoref*{#1}\ \nameref*{#1}}}

\begin{document}

\uksrcstandardfrontmatter
% Advanced/rare option if you specifically need a second metadata table:
% \ResetDocumentInfoRows
% \AddDocumentInfoItem{Document ID}{\DocumentID}
% \AddDocumentInfoItem{Document type}{\DocumentType}
% \AddDocumentInfoItem{Revision}{\DocumentRevision}
% \AddDocumentInfoItem{Classification}{\DocumentClassification}
% \AddDocumentInfoItem{Status}{\DocumentStatus}
% \AddDocumentInfoItem{Tags}{\DocumentTags}
% \AddDocumentInfoItem{Date}{\DocumentDate}
% \section*{Document information}
% \DocumentInfoTable
%
% Optional front matter if needed:
\uksrclistofabbreviations

\section{Introduction}

State the purpose of the memo, the intended audience, and the scope. Cite literature where relevant, for example \citet{1933Natur.132...66J}. Replace or add to this reference with a paper that is relevant to the memo. This sentence includes an example footnote\footnote{Example footnote text.} to show how notes are styled in the template.

\section{Context and Scope}

Use this optional section only if separate background or scope definition is useful.

\section{Main Discussion}

Use this section for the core material of the memo. Replace the example subsections with headings that fit the actual content.

\subsection{Topic 1}

Add the first main topic here.

\subsection{Topic 2}

Add the second main topic here.

\section{Conclusions and Recommendations}

Summarise the key outcome of the memo and any recommendations or decisions.

% Optional if there are unresolved points or follow-up actions to track.
% \section{Open Issues and Next Steps}
% Track unresolved points, dependencies, and next steps.
% Authors may optionally organise the main discussion around:
% - options considered
% - decision or recommendation
% - risks and limitations
%
% Optional back matter if needed:
% \uksrclistoffigures
% \uksrclistoftables

\bibliographystyle{plainnat}
\bibliography{uksrc-memo-template}

% Backup option for very short memos:
% \begin{thebibliography}{}
% \bibitem[Example(2026)]{example-ref}
% Example reference entry.
% \end{thebibliography}

\end{document}
